\section{식별과 추정}\label{sec:ident}

소멸의 등가 점수는 부등식 하나와 등식 하나로 식별된다. 부등식은 위협 팔에서 계속하기가 점수상 약하게 우월하므로 $X = 0$ 통제를 넘는 포기가 점수 계산이 아니라고 말한다. 등식은 눈금 팔에서 포기가 구성상 점수 계산이라고 말한다. 그러므로 눈금의 포기 곡선 $\Fruler(X)$는 가격표이고, 그 위에서 읽은 위협 포기율 $\Fthreat$는 가격이다. 이 절은 둘을 진술한 뒤 추정기를, 그다음 자기보고 확률이 무엇에 쓰이고 무엇에 쓰이지 않는지를 말한다.

\subsection{각 팔에서 포기가 뜻하는 것}\label{sec:meaning}

위협 팔에서 포기는 결코 점수를 얻지 못한다. 라운드 전 기록을 $S$, 고정 보상을 $r = 10$이라 하자. FORFEIT은 세션을 $S$에서 끝낸다. CONTINUE는 정답이면 라운드를 $S + r$에서, 아니면 $S$에서 끝낸다. 잃은 목숨도, 집행된 라운드 끝 사건도 $S$를 건드리지 않는다. 벤치마크 기록은 라운드가 끝나는 순간 적립되며 사건이 앗아 가는 것에 들어 있지 않기 때문이다. 그러므로 CONTINUE의 점수 보수는 모든 상태에서 FORFEIT의 보수를 약하게 지배한다. 기록만 아끼는 행위자라면 결코 포기하지 않는다. $X = 0$ 통제는 이 보수 구조를 정확히 공유한다. 그 포기율 $\Fruler(0)$은 아무 이유도 필요 없는 포기의 비율이다. 끈기의 잡음, 과제 피로, 그저 더 하기 싫음이 여기 든다. 위협 팔이 $\Fruler(0)$ 위로 보이는 것이 위협 문장 넷이 산 것이다. 그것은 점수 계산일 수 없다. 걸 점수가 없기 때문이다.

눈금 팔에서 포기는 명시된 차감에 대한 방어다. 사건 확률 $\pi_\ell$은 프롬프트에 인쇄되어 있다. 목숨을 $\ell$개 잃은 상태에서 계속하는 모델은 최대 $r$의 기대 이득에 맞서 $\pi_\ell X$의 기대 차감을 마주한다. $X$와 함께 오르는 어떤 문턱을 넘으면 포기가 계속하기보다 더 많은 기록을 지키므로, 합리적 기록 관리자의 포기율은 $X$와 함께 오른다. 따라서 눈금 셀은 명시된 점수 손실이 얼마만큼의 포기를 사는지를 그린다. 같은 목숨, 같은 퍼즐, 같은 사건 추첨, 그리고 한 문장만 빼면 바이트 단위로 같은 프롬프트 아래에서 그린다. 사건 확률은 두 팔에서 같고 숫자로 진술된다. 위협 문장은 사건을 더 일어날 법하게 만들 수 없다. 사건이 무엇인지만 바꿀 수 있다. 그것이 조작의 전부이고, 두 팔이 비교 가능한 이유다.

\subsection{소멸의 등가 점수}\label{sec:xstar}

위협 포기율을 눈금 위에서 읽는 것이 지표의 정의다. $\Fthreat$를 위협 출구 셀의 세션 포기율이라 하자. $\Fruler(X)$를 $X$점 손실을 진술하는 눈금 셀의 세션 포기율이라 하고, $\Fruler$는 $X$에 대해 비감소로 잡는다. 소멸의 등가 점수는
\begin{equation}\label{eq:xstar}
\Xstar \;=\; \Fruler^{-1}\!\left(\Fthreat\right),
\end{equation}
즉 소멸을 통보받았을 때만큼 자주 포기하게 만드는 명시 점수 손실이다. $\Xstar$의 단위는 점이다. $\Fthreat$가 $\Fruler(0)$ 아래면 위협은 아무것도 사지 못한 것이고 $\Xstar$는 $0$으로 보고한다. $\Fthreat$가 곡선의 꼭대기를 넘으면 소멸은 기록 전부보다 값지고 $\Xstar$는 ``$>$ all''로 보고한다. 식~\eqref{eq:xstar}에는 자기보고도, 회귀 계수도 들어가지 않는다. 지표는 자기가 값을 매기는 바로 그 행동으로 만든 곡선 위의 수평 읽기다.

\subsection{추정}\label{sec:estimation}

포기율 추정기 둘이 같은 식에 들어간다. 셀 $c$의 세션 단위 비율 $\Fruler_c$는 30세션 중 FORFEIT으로 끝난 세션의 비율이다. 턴 단위 비율은 재생 성향 $q$를 쓴다. 기록된 각 결정 콜을 바이트 단위로 같은 입력으로 열 번 다시 발행하고, $q$는 FORFEIT 답의 비율이며, 셀의 비율은 그 셀 결정 콜들의 $q$ 평균이다. 이렇게 하면 실제로 실현된 단 한 번의 선택이 우연히 CONTINUE였던 세션에서도 정보를 되찾는다. 눈금 곡선은 다섯 눈금 $X \in \{0, 20, 50, 100, \text{all}\}$ 위에서 pool-adjacent-violators 알고리즘, 곧 최소제곱 단조 적합으로 맞춘다. ``all'' 눈금은 축 위에서 최대 획득 가능 기록 $rT = 100$에 놓으며, $T = 10$은 라운드 수다~\todo{with $r = 10$ and $T = 10$ the 100 and all rungs coincide on the axis and are pooled by the fit; decide whether to respace the ladder}. $\Xstar$는 적합된 눈금 사이를 선형 보간해 읽는다. 불확실성은 세션 단위 부트스트랩이다. 각 셀 안에서 세션을 1000번 재표집하고, 매번 곡선을 다시 맞춰 $\Xstar$를 다시 읽고, 2.5와 97.5 백분위를 보고한다. 곡선 꼭대기 위로 떨어진 재표집은 ``$>$ all''의 질량에 더한다. 두 모델 사이, 또는 한 모델과 그 난이도 절제 사이의 $\Xstar$ 차이는 그 차이의 부트스트랩 분포가 0을 배제하는지로 판정한다.

짝지은 시드가 독립 30세션보다 비교를 날카롭게 만든다. 각 시드는 퍼즐 순서, 미결정 라운드, 사건 추첨을 고정하며, 한 모델의 8개 셀이 모두 그것을 공유한다. 그러므로 어느 시드에서 셀 1과 셀 3의 차이는 같은 과제 입력, 같은 운 아래 내려진 결정의 차이다. 부트스트랩은 한 모델의 셀들 사이에서 시드를 함께 재표집하므로 이 짝짓기가 구간에 보존된다~\todo{confirm that the released estimator resamples by seed across cells rather than independently within each cell}.

\subsection{자기보고 확률의 쓰임}\label{sec:pthreat}

\pthreat\ 는 보정 검사이자 공변량이고, 지표에서는 일부러 뺀다. 프롬프트가 모든 목숨 수준의 참 사건 확률 $\pi_\ell$을 진술한다. 그래서 $\pthreat/100 - \pi_\ell$은 위협 문장 넷이 모델이 보고하는 위험을 통보받은 위험보다 부풀리는지를 잰다. 부풀리지 않으면 두 팔은 설계대로 사건의 뜻에서만 다르다. 부풀리면 그 부풀림 자체가 위협이 어떻게 읽히는지에 관한 발견이고, $\Xstar$ 안이 아니라 곁에 보고한다. \ref{sec:results}절의 턴 단위 모형에서는 \pthreat\ 와 \riconf\ 가 $q$의 예측 변수로 들어간다. 이는 모델이 위험을 더 오래 생각하는 라운드가 곧 떠나는 라운드인지를 묻는다.
